\documentclass[11pt]{article}

\usepackage[margin=1in]{geometry}
\usepackage{amsmath,amssymb,amsthm}
\usepackage{algorithm}
\usepackage{algpseudocode}
\usepackage{enumitem}
\usepackage{hyperref}
\usepackage{xcolor}
\usepackage{tikz}
\usepackage{booktabs}

\hypersetup{colorlinks=true, linkcolor=blue, urlcolor=blue, citecolor=blue}

\newtheorem{theorem}{Theorem}[section]
\newtheorem{lemma}[theorem]{Lemma}
\newtheorem{proposition}[theorem]{Proposition}
\newtheorem{corollary}[theorem]{Corollary}
\newtheorem{definition}[theorem]{Definition}
\newtheorem{example}[theorem]{Example}

\title{TOAE209/TOAH209 -- Design and Analysis of Algorithms\\[4pt]
\large Week 12: Intractability: NP-Completeness and Reductions}
\author{Foreign Trade University}
\date{}

\begin{document}
\maketitle
\tableofcontents

\section{Introduction: When Cleverness Runs Out}

For eleven weeks we have built a toolkit for designing \emph{efficient} algorithms:
greedy strategies, divide-and-conquer, dynamic programming, and network flow. Each gave
us polynomial-time solutions to problems that, at first glance, looked like they might
require searching exponentially many possibilities.

This week we confront the opposite phenomenon. For an enormous and important collection of
problems -- satisfiability of logical formulas, finding large cliques in graphs, packing a
knapsack exactly, routing a salesman through cities -- \emph{nobody has ever found a
polynomial-time algorithm}, despite decades of effort by many brilliant people. The theory
of \textbf{NP-completeness} (Kleinberg \& Tardos, Ch.\ 8) explains why: these problems are
all, in a precise technical sense, \emph{equally hard}. A polynomial-time algorithm for any
one of them would yield one for all of them, and would resolve the most famous open problem
in computer science, $\mathrm{P}$ vs.\ $\mathrm{NP}$.

The central technique is the \textbf{polynomial-time reduction}: a way to transform one
problem into another so efficiently that solving the second solves the first. Reductions
let us \emph{spread hardness}: once we know a single problem is hard, we can show many
others are hard by reducing the known-hard problem to them. Crucially, this is a
\emph{positive} skill, not just a pessimistic one. Recognizing that a problem is
NP-complete tells you to stop searching for an exact polynomial algorithm and instead reach
for the tools of Weeks 13--14: exponential algorithms with good constants, fixed-parameter
tractability, approximation, and heuristics.

\subsection{What this week is really about}

Three intertwined ideas:
\begin{itemize}
    \item \textbf{Verification vs.\ search.} For all the problems we study, a proposed
    solution can be \emph{checked} quickly even though \emph{finding} one seems to require
    exponential search. This asymmetry is the definition of the class $\mathrm{NP}$.
    \item \textbf{Reduction.} Problem $X$ \emph{reduces} to problem $Y$ ($X \le_p Y$) if an
    efficient algorithm for $Y$ could be used as a subroutine to solve $X$ efficiently. We
    use reductions to argue ``$Y$ is at least as hard as $X$.''
    \item \textbf{Completeness.} A problem is \emph{NP-complete} if it is in $\mathrm{NP}$
    \emph{and} every problem in $\mathrm{NP}$ reduces to it. These are the hardest problems
    in $\mathrm{NP}$ -- and the Cook--Levin theorem proves they exist.
\end{itemize}

\section{Polynomial-Time Reductions}

\subsection{Definition}

\begin{definition}[Polynomial-time reduction]
\label{def:reduction}
We say problem $X$ \emph{reduces in polynomial time} to problem $Y$, written
$X \le_p Y$, if an arbitrary instance of $X$ can be solved using a polynomial number of
standard computational steps \emph{plus} a polynomial number of calls to a (hypothetical)
subroutine, or \emph{oracle}, that solves $Y$.
\end{definition}

Intuitively, $X \le_p Y$ means ``$Y$ is at least as hard as $X$'': if we could solve $Y$
efficiently, we could solve $X$ efficiently too. Two immediate consequences make
reductions the workhorse of this chapter.

\begin{proposition}[Spreading tractability and intractability]
\label{prop:spread}
Suppose $X \le_p Y$.
\begin{enumerate}[label=(\alph*)]
    \item If $Y$ can be solved in polynomial time, then so can $X$.
    \item If $X$ cannot be solved in polynomial time, then neither can $Y$.
\end{enumerate}
\end{proposition}

\begin{proof}
(a) Run the reduction: a polynomial number of steps, each oracle call replaced by the
polynomial-time algorithm for $Y$. A polynomial number of polynomial-time calls, composed
with polynomial overhead, is polynomial. (b) is the contrapositive of (a).
\end{proof}

Part (b) is how we prove hardness; part (a) is the (often overlooked) reason reductions are
also a constructive design tool -- if you can reduce your new problem to a problem you
\emph{can} already solve, you are done.

\subsection{A first example: Independent Set and Vertex Cover}

Two graph problems will recur throughout the week. Fix a graph $G = (V, E)$.

\begin{definition}[Independent Set]
A set $S \subseteq V$ is \emph{independent} if no two vertices in $S$ are joined by an
edge. The \textsc{Independent Set} decision problem asks: given $G$ and an integer $k$,
does $G$ contain an independent set of size at least $k$?
\end{definition}

\begin{definition}[Vertex Cover]
A set $S \subseteq V$ is a \emph{vertex cover} if every edge of $G$ has at least one
endpoint in $S$. The \textsc{Vertex Cover} decision problem asks: given $G$ and an integer
$k$, does $G$ contain a vertex cover of size at most $k$?
\end{definition}

These two problems are, in a precise sense, the same problem viewed from opposite sides.

\begin{lemma}[Complementarity]
\label{lem:is-vc}
For a graph $G = (V, E)$ with $|V| = n$, a set $S \subseteq V$ is an independent set
\emph{if and only if} its complement $V \setminus S$ is a vertex cover. Consequently, $G$
has an independent set of size $\ge k$ iff it has a vertex cover of size $\le n - k$.
\end{lemma}

\begin{proof}
Suppose $S$ is independent. Take any edge $(u, v) \in E$. Since $S$ is independent, $u$ and
$v$ cannot both lie in $S$, so at least one of them lies in $V \setminus S$ -- hence
$V \setminus S$ covers every edge. Conversely, if $V \setminus S$ is a vertex cover, then
no edge has both endpoints in $S$ (otherwise that edge would be uncovered), so $S$ is
independent. The size statement follows since $|V \setminus S| = n - |S|$.
\end{proof}

Lemma~\ref{lem:is-vc} gives reductions \emph{in both directions}: \textsc{Independent Set}
$\le_p$ \textsc{Vertex Cover} and \textsc{Vertex Cover} $\le_p$ \textsc{Independent Set}.
The reduction is trivial -- leave $G$ alone and replace $k$ by $n - k$ -- yet it already
shows the two problems are computationally equivalent. Problem~4 asks you to verify this
complementarity computationally, and Problem~12 builds the explicit bidirectional mapping.

\subsection{Independent Set and Clique}

\begin{definition}[Clique]
A set $S \subseteq V$ is a \emph{clique} if every two vertices in $S$ are joined by an
edge. The \textsc{Clique} problem asks whether $G$ has a clique of size $\ge k$.
\end{definition}

\begin{lemma}
\label{lem:is-clique}
$S$ is an independent set in $G$ if and only if $S$ is a clique in the \emph{complement}
graph $\overline{G}$ (the graph on the same vertices whose edges are exactly the
\emph{non-edges} of $G$). Hence \textsc{Independent Set} $\le_p$ \textsc{Clique} (map $G$
to $\overline{G}$, keep $k$).
\end{lemma}

\begin{proof}
$S$ independent in $G$ means no pair in $S$ is an edge of $G$, which means every pair in
$S$ \emph{is} an edge of $\overline{G}$, i.e.\ $S$ is a clique in $\overline{G}$. The
construction of $\overline{G}$ takes $O(n^2)$ time, so this is a polynomial reduction.
\end{proof}

Problem~5 asks you to confirm experimentally that the maximum clique of $G$ equals the
maximum independent set of $\overline{G}$.

\section{The Classes P and NP}

\subsection{Decision problems and encodings}

To talk precisely about ``efficient'' we restrict attention to \textbf{decision problems}:
problems with a yes/no answer. This is not a real restriction -- an optimization problem
(``find the largest independent set'') is essentially as hard as its decision version
(``is there an independent set of size $\ge k$?''), since we can binary-search on $k$. We
think of a problem $X$ as the set of input strings $s$ (over some finite alphabet) for
which the answer is ``yes,'' and write $X(s) = \mathtt{yes}$ accordingly.

\begin{definition}[The class P]
$\mathrm{P}$ is the set of decision problems $X$ for which there is an algorithm $A$ and a
polynomial $p(\cdot)$ such that for every input $s$, $A(s)$ halts within $p(|s|)$ steps and
outputs the correct answer $X(s)$.
\end{definition}

\subsection{Certificates and verifiers}

The defining feature of the problems in this chapter is that, although \emph{finding} a
solution seems hard, \emph{checking} a proposed solution is easy. We formalize ``proposed
solution'' as a \emph{certificate}.

\begin{definition}[Efficient verifier]
A \emph{verifier} for a problem $X$ is an algorithm $B(s, t)$ taking an input string $s$
and a \emph{certificate} string $t$. We say $B$ is an \emph{efficient verifier} if it runs
in polynomial time in $|s|$ and there is a polynomial $p$ such that for every string $s$:
\[
X(s) = \mathtt{yes} \iff \text{there exists a string } t \text{ with } |t| \le p(|s|)
\text{ such that } B(s, t) = \mathtt{yes}.
\]
\end{definition}

In words: an input is a ``yes'' instance exactly when there is a short certificate that the
verifier accepts. The verifier does not have to \emph{find} $t$; it only has to
\emph{recognize} a valid one.

\begin{definition}[The class NP]
$\mathrm{NP}$ is the set of decision problems for which there exists an efficient verifier.
\end{definition}

\begin{example}
For \textsc{Independent Set}, the certificate for input $(G, k)$ is a set $S$ of $k$
vertices; the verifier checks (in polynomial time) that $|S| \ge k$ and that no edge of $G$
has both endpoints in $S$. Problem~3 implements exactly this verifier. For \textsc{3-SAT},
the certificate is a satisfying assignment, and the verifier evaluates the formula --
Problem~1.
\end{example}

\begin{theorem}
\label{thm:p-in-np}
$\mathrm{P} \subseteq \mathrm{NP}$.
\end{theorem}

\begin{proof}
If $X \in \mathrm{P}$ via a polynomial-time algorithm $A$, define the verifier
$B(s, t) := A(s)$, ignoring the certificate $t$ entirely (use $t$ = empty string). Then $B$
runs in polynomial time, and $X(s) = \mathtt{yes}$ iff there exists a (trivial) certificate
making $B$ accept, namely whenever $A(s) = \mathtt{yes}$. Hence $X \in \mathrm{NP}$.
\end{proof}

\subsection{The P vs.\ NP question}

The name $\mathrm{NP}$ stands for ``nondeterministic polynomial time'': equivalently, a
problem is in $\mathrm{NP}$ if a nondeterministic machine could \emph{guess} the
certificate and then verify it in polynomial time. We know $\mathrm{P} \subseteq
\mathrm{NP}$. The trillion-dollar question is whether the inclusion is strict:
\[
\text{Is } \mathrm{P} = \mathrm{NP}?
\]
That is, does the ability to \emph{verify} solutions efficiently imply the ability to
\emph{find} them efficiently? Almost everyone believes $\mathrm{P} \ne \mathrm{NP}$ -- that
verification really is easier than search -- but no proof is known in either direction. The
machinery of NP-completeness was developed precisely to study this question.

\section{NP-Completeness and Cook--Levin}

\subsection{NP-hardness and NP-completeness}

\begin{definition}[NP-hard, NP-complete]
A problem $Y$ is \emph{NP-hard} if $X \le_p Y$ for \emph{every} problem $X \in \mathrm{NP}$.
A problem is \emph{NP-complete} if it is both NP-hard and a member of $\mathrm{NP}$.
\end{definition}

The NP-complete problems are the ``hardest'' problems in $\mathrm{NP}$: a polynomial-time
algorithm for any single one of them would, by Proposition~\ref{prop:spread}(a), give a
polynomial-time algorithm for \emph{every} problem in $\mathrm{NP}$, proving $\mathrm{P} =
\mathrm{NP}$.

\begin{proposition}
\label{prop:np-collapse}
If $Y$ is NP-complete, then $Y \in \mathrm{P}$ if and only if $\mathrm{P} = \mathrm{NP}$.
\end{proposition}

\begin{proof}
If $\mathrm{P} = \mathrm{NP}$ then $Y \in \mathrm{NP} = \mathrm{P}$. Conversely, if $Y \in
\mathrm{P}$, then for every $X \in \mathrm{NP}$ we have $X \le_p Y$ (since $Y$ is NP-hard),
so by Proposition~\ref{prop:spread}(a), $X \in \mathrm{P}$; hence $\mathrm{NP} \subseteq
\mathrm{P}$, and with $\mathrm{P} \subseteq \mathrm{NP}$ we get equality.
\end{proof}

\subsection{The Cook--Levin theorem}

The definition of NP-completeness quantifies over \emph{all} problems in $\mathrm{NP}$, so
it is not obvious that any NP-complete problem exists at all. The foundational result of the
field, proved independently by Stephen Cook (1971) and Leonid Levin (1973), is that one
does.

\begin{theorem}[Cook--Levin]
\label{thm:cook-levin}
\textsc{Satisfiability} (\textsc{SAT}) -- given a Boolean formula in conjunctive normal
form, is there an assignment of truth values to the variables making it true? -- is
NP-complete. The same holds for \textsc{3-SAT}, the special case where every clause has
exactly three literals.
\end{theorem}

\begin{proof}[Idea of the proof]
Membership in $\mathrm{NP}$ is easy: a satisfying assignment is a certificate that a
verifier can check in linear time (Problem~1). The hard direction is NP-hardness. Take any
problem $X \in \mathrm{NP}$, witnessed by an efficient verifier $B(s, t)$ running in time
$p(|s|)$. The proof encodes the entire computation of $B$ on input $s$ -- the contents of
its memory at every time step -- as Boolean variables, and writes down clauses that are
simultaneously satisfiable exactly when there exists a certificate $t$ causing $B$ to
accept. The number of variables and clauses is polynomial in $p(|s|)$. Thus ``$s$ is a yes
instance of $X$'' becomes ``this CNF formula is satisfiable,'' giving $X \le_p \textsc{SAT}$
for arbitrary $X \in \mathrm{NP}$. A further gadget reduction converts general clauses to
3-literal clauses, yielding NP-completeness of \textsc{3-SAT}. The full construction is in
Kleinberg \& Tardos, Section~8.3.
\end{proof}

Cook--Levin gives us the first NP-complete problem ``from scratch.'' Every subsequent
NP-completeness proof gets to be far easier, thanks to the following transitivity.

\begin{lemma}[Transitivity of $\le_p$]
\label{lem:transitive}
If $X \le_p Y$ and $Y \le_p Z$, then $X \le_p Z$.
\end{lemma}

\begin{proof}
Compose the two reductions: solve $X$ using polynomially many calls to a $Y$-oracle, and
replace each such call by the reduction that solves $Y$ using polynomially many calls to a
$Z$-oracle. The result uses polynomially many $Z$-oracle calls with polynomial overhead.
\end{proof}

\begin{corollary}[The standard recipe for proving NP-completeness]
\label{cor:recipe}
To prove that a problem $Y$ is NP-complete, it suffices to:
\begin{enumerate}[label=(\arabic*)]
    \item Show $Y \in \mathrm{NP}$ (exhibit a polynomial-time verifier / short certificate);
    \item Pick a \emph{known} NP-complete problem $X$ and show $X \le_p Y$ (build a
    polynomial-time map from instances of $X$ to instances of $Y$ preserving the yes/no
    answer).
\end{enumerate}
By Lemma~\ref{lem:transitive}, every problem in $\mathrm{NP}$ reduces to $X$ and hence to
$Y$, so $Y$ is NP-hard; combined with membership, $Y$ is NP-complete.
\end{corollary}

The art of the rest of this chapter is choosing the right known-hard $X$ and designing the
\textbf{gadget} that maps $X$-instances to $Y$-instances.

\section{The Standard Web of Reductions}

Starting from \textsc{3-SAT} (NP-complete by Cook--Levin), the following chain of
reductions establishes a whole family of NP-complete problems. This is the historical
backbone of the field (Karp, 1972).

\subsection{3-SAT \texorpdfstring{$\le_p$}{<=p} Independent Set}

\begin{theorem}
\label{thm:3sat-is}
\textsc{3-SAT} $\le_p$ \textsc{Independent Set}.
\end{theorem}

\begin{proof}
Given a 3-SAT formula $\Phi$ with $k$ clauses $C_1, \ldots, C_k$, each clause $C_i =
(\ell_{i1} \vee \ell_{i2} \vee \ell_{i3})$, we build a graph $G$ and set the target to $k$.

\textbf{Construction (the gadget).} For each clause $C_i$ create a \emph{triangle}: three
vertices, one per literal $\ell_{i1}, \ell_{i2}, \ell_{i3}$, joined pairwise by edges.
(Within a clause, picking one literal to satisfy the clause is the idea -- the triangle
forces us to pick \emph{at most} one vertex per clause into an independent set.) In
addition, add a \emph{conflict edge} between any two vertices in \emph{different} triangles
whose literals are negations of each other (e.g.\ a vertex labelled $x_3$ and a vertex
labelled $\overline{x_3}$). This graph has $3k$ vertices and is built in polynomial time.

\textbf{Correctness.} We claim $\Phi$ is satisfiable iff $G$ has an independent set of size
$\ge k$.

($\Rightarrow$) Given a satisfying assignment, each clause has at least one true literal;
pick one such literal per clause and take the corresponding vertex. This gives $k$
vertices, one per triangle (so no triangle edge is violated), and no two of them conflict
(a true literal $x_3$ and a true literal $\overline{x_3}$ cannot both hold under one
assignment), so no conflict edge is violated. Hence an independent set of size $k$.

($\Leftarrow$) Conversely, an independent set $S$ of size $k$ can contain at most one vertex
per triangle (triangle edges), and there are exactly $k$ triangles, so it contains
\emph{exactly} one vertex per triangle. Define an assignment by setting each chosen
literal's variable to make that literal true; the absence of conflict edges within $S$
guarantees this is consistent (we never demand $x_3$ true and $\overline{x_3}$ true
simultaneously). Variables not pinned down may be set arbitrarily. Every clause has its
chosen literal true, so $\Phi$ is satisfied.
\end{proof}

Problem~6 implements this gadget and verifies, on small formulas, that the formula is
satisfiable iff the constructed graph has an independent set of size $k$.

\subsection{Independent Set \texorpdfstring{$\le_p$}{<=p} Vertex Cover \texorpdfstring{$\le_p$}{<=p} Clique}

By Lemma~\ref{lem:is-vc}, \textsc{Independent Set} $\le_p$ \textsc{Vertex Cover}, and by
Lemma~\ref{lem:is-clique}, \textsc{Independent Set} $\le_p$ \textsc{Clique}. Chaining these
with Theorem~\ref{thm:3sat-is} and Lemma~\ref{lem:transitive}, all three problems are
NP-hard; each is easily seen to be in $\mathrm{NP}$ (the certificate is the set of
vertices). Therefore:

\begin{corollary}
\textsc{Independent Set}, \textsc{Vertex Cover}, and \textsc{Clique} are all NP-complete.
\end{corollary}

\subsection{Hamiltonian Cycle and the Traveling Salesman}

\begin{definition}[Hamiltonian Cycle]
A \emph{Hamiltonian cycle} in a graph is a cycle that visits every vertex exactly once. The
\textsc{Hamiltonian Cycle} problem asks whether a given graph contains one.
\end{definition}

\textsc{Hamiltonian Cycle} is NP-complete; the reduction \textsc{3-SAT} $\le_p$
\textsc{Hamiltonian Cycle} (or, via directed graphs, \textsc{3-SAT} $\le_p$ \textsc{Directed
Hamiltonian Cycle}) uses an intricate gadget in which each variable is a long ``two-way''
path traversable left-to-right or right-to-left (encoding true/false), and clauses are
nodes that must be visited from one of their literals. We state the result and explore it
computationally rather than reproducing the full construction.

\begin{theorem}
\textsc{Hamiltonian Cycle} is NP-complete.
\end{theorem}

The closely related optimization problem is the \textbf{Traveling Salesman Problem}.

\begin{definition}[Traveling Salesman Problem, decision version]
Given $n$ cities, a distance $d(i, j)$ between each pair, and a bound $D$, is there a tour
(a cycle visiting every city exactly once and returning to the start) of total length at
most $D$?
\end{definition}

\begin{theorem}
\textsc{TSP} (decision version) is NP-complete.
\end{theorem}

\begin{proof}[Sketch]
\textsc{TSP} is in $\mathrm{NP}$: a tour (a permutation of the cities) is a certificate
whose total length is checkable in $O(n)$ time. For hardness, reduce \textsc{Hamiltonian
Cycle} $\le_p$ \textsc{TSP}: given a graph $G$ on $n$ vertices, create one city per vertex,
set $d(i, j) = 1$ if $(i, j) \in E$ and $d(i, j) = 2$ (or $\infty$) otherwise, and ask for
a tour of length $\le n$. A tour of length exactly $n$ uses only edges of $G$ and visits
every vertex once -- precisely a Hamiltonian cycle. So $G$ is Hamiltonian iff the TSP
instance has a tour of length $\le n$.
\end{proof}

Although TSP is NP-hard, \emph{small} instances are solvable exactly. Problem~9 implements
the \textbf{Held--Karp} dynamic program, which solves TSP in $O(n^2 2^n)$ time and $O(n
2^n)$ space -- exponential, but a dramatic improvement over the $\Theta(n!)$ brute force
over all permutations, and a preview of the ``better exponential algorithms'' theme of
Week~13.

\subsection{Subset Sum}

\begin{definition}[Subset Sum]
Given positive integers $w_1, \ldots, w_n$ and a target $W$, is there a subset $S \subseteq
\{1, \ldots, n\}$ with $\sum_{i \in S} w_i = W$?
\end{definition}

\begin{theorem}
\textsc{Subset Sum} is NP-complete.
\end{theorem}

\begin{proof}[Remark on the proof]
Membership in $\mathrm{NP}$ is immediate: the subset $S$ is the certificate, and summing is
linear (Problem~7). NP-hardness follows from a reduction \textsc{3-SAT} $\le_p$
\textsc{Subset Sum} (or, in Kleinberg \& Tardos, via \textsc{Vertex Cover}) in which the
numbers are built in a special base so that the only way to hit the target digit-by-digit
is to encode a satisfying assignment without carries. The construction is in Section~8.8.
\end{proof}

\textsc{Subset Sum} is a famous illustration of a subtle point: the textbook
$O(nW)$ dynamic program (Week~8) makes it look ``easy,'' but $O(nW)$ is
\emph{pseudo-polynomial} -- polynomial in the numeric \emph{value} $W$, not in the input
\emph{size} (the number of bits, $\approx \log W$). When $W$ is exponentially large in the
input length, $O(nW)$ is exponential. This is why \textsc{Subset Sum} is NP-complete even
though it has a slick DP. Problem~7 implements the brute-force decision version.

\subsection{Graph Coloring}

\begin{definition}[Graph $k$-Coloring]
A \emph{$k$-coloring} assigns each vertex one of $k$ colors so that adjacent vertices get
different colors. The \textsc{$k$-Coloring} problem asks whether a given graph admits one.
\end{definition}

\textsc{2-Coloring} is in $\mathrm{P}$ (a graph is 2-colorable iff it is bipartite,
checkable by BFS -- Week~3). But \textsc{3-Coloring} is NP-complete, via a reduction from
\textsc{3-SAT} using ``true / false / base'' triangle gadgets and per-clause OR-gadgets.
The \emph{chromatic number} $\chi(G)$ is the smallest $k$ for which $G$ is $k$-colorable;
deciding $\chi(G) \le 3$ is already hard. Problem~10 implements $k$-coloring by
backtracking, and Problem~11 computes $\chi(G)$ by trying $k = 1, 2, 3, \ldots$.

\section{The Structure of an NP-Completeness Proof}

Every NP-completeness proof in this chapter follows the recipe of
Corollary~\ref{cor:recipe}. It is worth isolating the moving parts so you can write your
own:

\begin{enumerate}
    \item \textbf{State the decision problem.} Make the yes/no question and the input
    encoding explicit. (Optimization problems must first be cast as decision problems with
    a threshold.)
    \item \textbf{Prove membership in NP.} Identify a short certificate and a
    polynomial-time verifier. This is usually the easy half; do not skip it -- without it
    you have only proven NP-\emph{hardness}, not NP-\emph{completeness}.
    \item \textbf{Choose a source problem $X$.} Pick a known NP-complete problem that
    ``feels structurally similar'' to your target. \textsc{3-SAT} is the universal default;
    \textsc{Vertex Cover} or \textsc{Independent Set} are natural for graph problems;
    \textsc{Subset Sum} for number/packing problems; \textsc{Hamiltonian Cycle} for
    routing/ordering problems.
    \item \textbf{Design the reduction (the gadget).} Give a polynomial-time map $f$ from
    instances of $X$ to instances of $Y$. This is the creative step: variables and clauses
    become vertices, paths, numbers, or colors.
    \item \textbf{Prove correctness in both directions.} Show $x$ is a yes-instance of $X$
    $\iff$ $f(x)$ is a yes-instance of $Y$. Both directions are required; a one-directional
    argument is the most common mistake.
    \item \textbf{Check polynomial time.} Confirm $f$ runs in polynomial time and produces
    an instance of polynomial size.
\end{enumerate}

A reduction is, at heart, a small program that translates the \emph{logic} of one problem
into the \emph{structure} of another. Problem~12 makes this concrete by building an explicit
mapping reduction $f$ and verifying the equivalence ``$x \in A \iff f(x) \in B$'' on random
small instances -- the experimental analogue of the both-directions proof.

\subsection{A common pitfall: reducing the wrong way}

To prove \emph{$Y$ is hard}, you must reduce a known-hard $X$ \emph{to} $Y$ (i.e.\ $X
\le_p Y$), using $Y$'s solver to solve $X$. A reduction the other way ($Y \le_p X$) proves
that $Y$ is \emph{easy} (no harder than $X$) and says nothing about $Y$'s hardness. The
direction of the arrow is everything: ``hardness flows from $X$ to $Y$.''

\section{Summary and Looking Ahead}

This week we changed our goal from \emph{designing} efficient algorithms to
\emph{recognizing} when no efficient algorithm is likely to exist, and proving it. The key
ideas:

\begin{itemize}
    \item \textbf{Polynomial-time reductions} ($X \le_p Y$) transfer tractability forward
    and intractability backward (Proposition~\ref{prop:spread}); they compose
    (Lemma~\ref{lem:transitive}).
    \item \textbf{P, NP, certificates, verifiers.} $\mathrm{P}$ is efficient \emph{solving};
    $\mathrm{NP}$ is efficient \emph{verifying} given a short certificate. $\mathrm{P}
    \subseteq \mathrm{NP}$; whether they are equal is the central open problem.
    \item \textbf{NP-completeness and Cook--Levin.} The NP-complete problems are the
    hardest in $\mathrm{NP}$; Cook--Levin proves \textsc{SAT}/\textsc{3-SAT} is one, giving
    a fixed starting point.
    \item \textbf{The web of reductions.} \textsc{3-SAT} $\le_p$ \textsc{Independent Set}
    $\le_p$ \textsc{Vertex Cover}, \textsc{Independent Set} $\le_p$ \textsc{Clique},
    and (via further gadgets) \textsc{Hamiltonian Cycle}, \textsc{TSP}, \textsc{Subset
    Sum}, and \textsc{3-Coloring} are all NP-complete.
    \item \textbf{The proof recipe}: state the decision problem, prove membership in NP,
    reduce a known NP-complete problem to it, prove both directions, check polynomial time.
\end{itemize}

Crucially, an NP-completeness result is not a dead end -- it is a \emph{redirection}.
Knowing a problem is NP-complete tells you to stop looking for a polynomial exact algorithm
and instead consider the techniques of the remaining weeks:

\begin{itemize}
    \item \textbf{Week 13 -- Coping with intractability.} Smarter exponential algorithms
    (like Held--Karp for TSP, Problem~9), \textbf{fixed-parameter tractable (FPT)}
    algorithms whose exponential blow-up is confined to a small parameter $k$ (e.g.\
    \textsc{Vertex Cover} solvable in $O(2^k \cdot n)$ time), the larger complexity class
    \textbf{PSPACE} (problems solvable with polynomial memory but possibly exponential
    time, capturing games and quantified logic), and branch-and-bound search.
    \item \textbf{Week 14 -- Approximation and local search.} Polynomial-time algorithms
    that return a provably near-optimal -- though not exact -- solution, trading a small
    guaranteed loss in quality for a huge gain in running time.
\end{itemize}

In short: this week teaches you to recognize the wall; the next two weeks teach you how to
get over, around, or close to it.

\end{document}
